\section{The Discounted Objective}
\label{sec:discounted}

The foundational paper's optimization loop selects actions by their immediate
effect:
\begin{equation}
\label{eq:myopic}
\pi^* = \arg\max_\pi \mathbb{E}[\Delta \volL(G) \mid \pi]
\end{equation}
This is the $\gamma = 0$ special case of a discounted multi-step objective.

\subsection{Definition}

\begin{definition}[Discounted Value Function]
\label{def:discounted}
For a discount factor $\gamma \in [0, 1)$ and a planning horizon $H \leq
    \infty$, the \emph{discounted value} of an action sequence
    $\boldsymbol{\pi} = (\pi_0, \pi_1, \ldots, \pi_{H-1})$ is:
\begin{equation}
\label{eq:discounted}
\Vdisc(\boldsymbol{\pi}) = \mathbb{E}\!\left[\sum_{t=0}^{H-1} \gamma^t
    \Delta \volL(G_t \mid \pi_t)\right]
\end{equation}
where $G_t$ is the capability poset state at time $t$, $\Delta \volL(G_t
    \mid \pi_t)$ is the volume change produced by action $\pi_t$ at time $t$,
    and the expectation is taken over the actor's world model
    $\WorldModel$ of future states.
\end{definition}

The discounted value function is a standard object in sequential decision
theory \citep{puterman1994markov, sutton2018reinforcement}. The
$\volL$-specific content is in what it sums: volume changes on the capability
poset, each satisfying the axioms M1--M6 established in
\citet{lasser2026poset}. The discount factor $\gamma$ controls the planning
horizon: $\gamma$ near 0 recovers the myopic objective; $\gamma$ near 1
extends planning to the far future. The choice of $\gamma$ is an
implementation parameter for the self-balancing transfer: different agents
may use different discount factors depending on their planning capacity and
the timescale of the problems they face. However, the anti-monopolar
property (Section~\ref{sec:substitution}) introduces theoretical constraints
on $\gamma$ under certain domination models.

\subsection{Self-Balancing Transfers to the Discounted Objective}

The following theorem is a \emph{consistency result}, not a sufficiency
result: it establishes that the per-step self-balancing property of
\citet{lasser2026gfm} composes correctly under discounted summation, so
the discounted objective inherits the same step-local guarantees as the
myopic objective. It does not, by itself, rule out the failure modes that
emerge specifically from multi-step planning --- in particular, sequences
where each step is locally self-balancing but the trajectory ends in
catastrophic contraction because the destructive consequence is deferred.
That class of failures requires the risk-capability machinery of
Section~\ref{sec:risk}, not the transfer theorem.

\begin{theorem}[Self-Balancing Under Discounting]
\label{thm:discounted_self_balancing}
Let $\volL$ satisfy axioms M1--M6 (Proposition~1 of
    \citet{lasser2026poset}).
    For any discount factor $\gamma \in [0, 1)$ and planning horizon $H$:
\begin{enumerate}
\item[\textbf{(a)}] \textbf{Per-step preservation.} At each time step $t$,
    Lemmas~1--3 of \citet{lasser2026gfm} apply to $\Delta \volL(G_t \mid
    \pi_t)$: destruction, coercion, and rigidity are each
    $\volL$-contracting at that step.
\item[\textbf{(b)}] \textbf{Aggregate preservation.} A sequence
    $\boldsymbol{\pi}$ in which every action is $\volL$-contracting has
    $\Vdisc(\boldsymbol{\pi}) < 0$. A sequence in which every action is
    $\volL$-expanding has $\Vdisc(\boldsymbol{\pi}) > 0$.
\item[\textbf{(c)}] \textbf{Mixed sequences.} A sequence containing both
    expanding and contracting actions has $\Vdisc$ determined by the
    $\gamma$-weighted balance. Temporary contraction followed by larger
    expansion is value-positive; sustained contraction is value-negative
    regardless of $\gamma$.
\item[\textbf{(d)}] \textbf{Structural balance.} The discounted objective
    penalizes sustained destructive permissiveness, sustained coercive
    restriction, and sustained rigidity. Temporary restriction
    (quarantine, preemptive containment) is value-positive when the
    discounted future expansion exceeds the current contraction.
\end{enumerate}
\end{theorem}

\begin{remark}[M4 across the two foundational papers]
\label{rem:m4_note}
The poset paper's axiom M4 uses \emph{poset-disjointness}, which is a
    stronger condition than the foundational paper's set-disjointness used
    in the original self-balancing proofs. The discounted transfer holds
    under the stronger M4 because on a poset with non-negative independent
    weights $\iw$, removing nodes from $\dc G$ can only reduce
    $\sum \iw$; the argument does not depend on whether M4-additivity
    holds for the specific partition used in the original proofs. Readers
    tracing Lemmas~1--3 across the two papers should keep this
    discrepancy in mind.
\end{remark}

\begin{proof}
(a) At each time step $t$, the poset state $G_t$ is a finite poset with
measure $\volL$ satisfying M1--M6 (the axioms depend on the poset structure,
which is maintained at each step by the insertion policy of
\citet{lasser2026poset}). Lemmas~1--3 are derived from M1--M6 alone and
therefore apply to $\Delta \volL(G_t \mid \pi_t)$ at every $t$.

(b) If $\Delta \volL(G_t \mid \pi_t) < 0$ for all $t$, then each term in
Equation~\eqref{eq:discounted} is negative and weighted by $\gamma^t > 0$. The
sum of negative terms with positive weights is negative:
$\Vdisc(\boldsymbol{\pi}) < 0$. The expanding case is symmetric.

(c) For a mixed sequence, $\Vdisc = \sum_t \gamma^t \Delta \volL(G_t \mid
\pi_t)$ is a weighted sum of positive and negative terms. The sign of
$\Vdisc$ depends on the magnitudes: a contraction of $-C$ at time $t_0$
followed by an expansion of $+E$ at time $t_1 > t_0$ gives a net
contribution of $-\gamma^{t_0} C + \gamma^{t_1} E$, which is positive when
$E > C \cdot \gamma^{t_0 - t_1}$. Since $\gamma < 1$, future expansions are
discounted relative to present contractions: the actor requires a larger
future payoff to justify current contraction, which is the standard
discounting tradeoff.

(d) We say a strategy exhibits \emph{sustained contraction} when
$\Delta \volL(G_t \mid \pi_t) \leq 0$ for all $t$, with strict inequality
for at least one step. By (b), any such strategy has $\Vdisc \leq 0$. More
generally, a strategy that contracts at every step beyond some finite prefix
has $\Vdisc$ determined by (c): the expanding prefix contributes a finite
positive sum while the contracting tail contributes a negative sum, and
the strategy is value-negative whenever the tail's magnitude exceeds
the prefix's. Temporary contraction (a finite number of contracting steps
followed by expanding steps) can be value-positive by (c). The
self-balancing property (Proposition~1 of \citet{lasser2026gfm}) holds at
each step by (a), and the discounted sum preserves the structural balance:
the objective penalizes sustained extremes and allows temporary restriction
only when the future expansion justifies it.
\end{proof}

\paragraph{What the theorem says.} The myopic objective treats every
contraction as bad. The discounted objective distinguishes between
\emph{sustained} contraction (always bad) and \emph{temporary} contraction
followed by expansion (potentially good). This is the formal basis for
quarantine, preemptive containment, and risk-capability restriction: the actor
can accept a small current $\volL$-loss to prevent a larger future loss, as
long as the discounted arithmetic works out. The self-balancing property is
not weakened; it is extended to handle temporal tradeoffs.

\paragraph{What the theorem does not say.} It does not specify the planning
algorithm (how the actor searches over action sequences), the world model
(how the actor predicts future $G_t$ states), or the discount factor (how the
actor balances present and future). These are implementation choices that
shape the quality of planning without affecting the structural guarantee. A
poor world model will produce poor plans, but even a poor plan is evaluated by
an objective that penalizes sustained contraction. The self-balancing property
is a property of the \emph{objective}, not of the \emph{planner}.\footnote{The
companion paper's leverage signal $\lambda(d)$
(Definition~9 of \citet{lasser2026poset}) is defined as an instantaneous
excess-growth-rate attributable to capability $d$ and is used by the
myopic score $S(\pi)$. Under the discounted objective the corresponding
quantity is the leverage integrated against the geometric weight
$\sum_{t \geq 0} \gamma^t \lambda(d_t)$, evaluated along the actor's
predicted trajectory. Because $\lambda$ is itself a $\volL$-rate, the
geometric sum is well-defined whenever $\lambda(d_t)$ is bounded along the
trajectory, and the discounted score reduces to the myopic
$\lambda$-weighted score in the $\gamma \to 0$ limit. The leverage
ranking signal is therefore consistent with the discounted objective in
the same sense that the myopic self-balancing property is consistent: each
per-step leverage estimate enters the discounted sum with the same
geometric weight as the corresponding $\Delta\volL$, so the leverage
landscape that directs the planner is the discounted version of the
companion paper's instantaneous landscape, not a separate object. A
formal restatement of the leverage estimator under discounting is
straightforward but is left to the implementation paper.}
